\section{Frontiers: Beyond a Single MDP}

% Shared colour palette for this lecture's hand-built diagrams.
% (No theme colour macros exist in the preamble; we use plain xcolor names
%  plus this small deck-local palette throughout the new TikZ. Names match
%  Day 9 for visual continuity across the two closing decks.)
\definecolor{sysblue}{RGB}{189,215,238}   % model / policy boxes
\definecolor{mlgreen}{RGB}{169,208,142}   % positive / shared structure
\definecolor{boxgreen}{RGB}{197,224,180}
\definecolor{boxblue}{RGB}{180,199,231}
\definecolor{boxred}{RGB}{244,177,177}    % warning / conflict / fragile
\definecolor{boxtan}{RGB}{248,203,173}    % task / human inputs
\definecolor{accentred}{RGB}{192,0,0}     % punchline callouts

% Source/citation footnote with no visible mark (suppresses the stray "0"
% that \footnote[0] leaves behind). Counter is restored so numbering is unaffected.
\newcommand{\srcnote}[1]{%
  \begingroup
  \renewcommand{\thefootnote}{}%
  \footnote{#1}%
  \addtocounter{footnote}{-1}%
  \endgroup}

% ---------------------------------------------------------------- Recap from Day 9
\begin{frame}{Where We Left Off: Day 9}
\begin{itemize}\itemsep0.5em
    \item Yesterday put RL in the \textbf{real world}: \textbf{RLHF} (learn the reward
          when you can't write it down), \textbf{multi-agent RL}, and \textbf{robotics}.
    \item Day 9 closed on a warning we pay off today:
          \begin{itemize}
              \item \textbf{multi-\emph{agent}} RL $=$ many agents sharing \emph{one}
                    environment (non-stationarity);
              \item \textbf{multi-\emph{task}} RL $=$ \emph{one} agent across \emph{many}
                    tasks. \empr{These are not the same thing.}
          \end{itemize}
    \item It also flagged two live frontiers: \textbf{emergent complexity} from self-play,
          and \textbf{verifiable rewards} (RLVR) for reasoning agents.
\end{itemize}
\vfill
\begin{center}
\itshape Every method in Days 1--9 optimised \emph{one} fixed MDP. Today: what happens when
we let go of that assumption?
\end{center}
\end{frame}

% ---------------------------------------------------------------- Today's map
\begin{frame}{Today: Four Ways RL Generalises Past One MDP}
\footnotesize Here a \emph{task} means one MDP: its own dynamics and reward (a specific
maze, a specific robot skill).\par\normalsize
\vspace{0.5em}
\begin{itemize}\itemsep0.5em
    \item \textbf{Meta-RL} --- \emph{learning to learn}: adapt to a \emph{new} task from a
          little experience (context \empr{inferred}).
    \item \textbf{Multi-task RL} --- one policy over a \emph{fixed set} of tasks (context
          \empr{given}); where negative transfer comes from.
    \item \textbf{Hierarchical RL} --- reasoning over \emph{skills}, not primitive actions:
          generalising across \emph{time-scales}.
    \item \textbf{Open problems} --- what still breaks, and a reflective close on where the
          field is going.
\end{itemize}
\vfill
\begin{center}
The thread: \empy{generalising beyond a single fixed problem} --- across tasks, to new
tasks, and across time.
\end{center}
\end{frame}

% ---------------------------------------------------------------- Learning outcomes
\begin{frame}{Learning Outcomes}
This is an \textbf{overview} day. By the end you should be able to:
\begin{itemize}\itemsep0.3em
    \item \textbf{define and place} meta-RL, multi-task RL, and hierarchical RL --- and say
          precisely how they differ;
    \item explain the core idea of meta-RL --- \empr{learning a \emph{learner}} $f_\theta$
          --- and how RNNs and transformers (\emph{in-context RL}) implement it;
    \item explain \textbf{why} each frontier is hard (exploration as a learned skill,
          negative transfer, discovering abstractions);
    \item name \textbf{representative methods} (RL$^2$ / Decision Transformer, PCGrad, the
          options framework) and articulate the main \textbf{open problems}, linking back to
          Days 1--9.
\end{itemize}
\vspace{0.4em}
\begin{center}\footnotesize\itshape
Goal: to \emph{recognise and situate} these frontiers, not to implement each from scratch.
\end{center}
\end{frame}

% ---------------------------------------------------------------- Course arc
\begin{frame}{The Whole Journey: Ten Days}
\centering
\resizebox{\textwidth}{!}{%
\begin{tikzpicture}[
    font=\footnotesize,
    node distance=0pt,
    stage/.style={draw, rounded corners, minimum height=1.05cm, align=center,
                  inner xsep=5pt, text width=2.1cm},
    ar/.style={-{Stealth[length=2mm]}, line width=1pt, black!55}]
  \node[stage, fill=sysblue]  (d1) {\textbf{1--2}\\ foundations\\ value / DQN};
  \node[stage, fill=sysblue, right=6mm of d1] (d3) {\textbf{3--4}\\ policy gradient\\ A2C / PPO};
  \node[stage, fill=boxgreen, right=6mm of d3] (d5) {\textbf{5--6}\\ continuous\\ DDPG / SAC};
  \node[stage, fill=boxgreen, right=6mm of d5] (d7) {\textbf{7}\\ exploration\\ \& reward / IRL};
  \node[stage, fill=boxtan, right=6mm of d7] (d8) {\textbf{8}\\ model-based\\ \& offline};
  \node[stage, fill=boxtan, right=6mm of d8] (d9) {\textbf{9}\\ real world\\ RLHF / robots};
  \node[stage, fill=accentred, text=white, right=6mm of d9] (d10) {\textbf{10}\\ \textbf{frontiers}\\ today};
  \draw[ar] (d1)--(d3); \draw[ar] (d3)--(d5); \draw[ar] (d5)--(d7);
  \draw[ar] (d7)--(d8); \draw[ar] (d8)--(d9); \draw[ar] (d9)--(d10);
\end{tikzpicture}}
\vfill
\begin{center}\footnotesize
From \emph{how to optimise one MDP} (Days 1--6) $\rightarrow$ \emph{where the reward and
data come from} (Days 7--8) $\rightarrow$ \emph{the real world} (Day 9) $\rightarrow$
\emph{past a single MDP entirely} (today).
\end{center}
\end{frame}
